\section{Compression from ports}\label{sec:compr}

Assume throughout this section that the characteristic admits ports.
To enumerate a contraction directly, one assigns a bit to every wire
and multiplies the entries of all tensors. The compressor replaces
this enumeration by a sum in which fewer bits are shared between
positions. Once those bits are fixed, each position can sum over its
own indices independently. We describe the expansion from
\cite[Theorems~\ext{G}{thm:compressor} and~\ext{G}{thm:conversion},
Lemma~\ext{G}{lem:descriptors}]{oaG} and use it to evaluate a closed
network in subexponential time. Section~\ref{sec:gap} will use the
separate bounds on shared and private work.

\subsection{An exact expansion with fewer shared bits}
The compressor applies to a network indexed by base vertices and lanes.
The base is a directed loop-free multigraph of fixed in- and out-degree
$d$, a power of two with $d\ge4$. Each pair $(u,x)$ of a base vertex
and a lane $x\in\{0,1\}^h$ is a \emph{position}. A base arc
$e:u\to v$ with mask $\tau_e$ induces an actual arc from $(u,x)$ to
$(v,x+\tau_e)$ in each lane, with addition bitwise modulo two. Each
actual arc carries a string of bits shared by its two endpoints. The
width may vary from lane to lane and may be zero. Write $P$ for the
number of positions, $D$ for the maximum width and $A$ for the
contraction of the factors at the positions. A factor may also
depend on parameters, including advice bits.

One compression block applies a fixed sequence of local tensor
identities. Each identity is a sum, and a \emph{history} records one
admissible summand at every application. Given $t\ge1$ blocks with
$b_0t\le h$, where $b_0$ is fixed, the final network in a history
$\lambda$ has bounded degree and at most $P$ positions. Let $z$ assign
bits to its arcs. At a final position $v$, let $y_v$ assign the bits
private to that position and let $K_{\lambda,v}$ be the product of its
local entries and equality tests. The compressor gives the identity
\[
 A=\sum_{\lambda}\sum_z
          \prod_v\left(\sum_{y_v}
             K_{\lambda,v}(z|_v,y_v)\right),
\]
where $z|_v$ denotes the strings on arcs incident to $v$. Each original
factor moves intact to one final position, with all of its arguments.
Three estimates determine the cost of implementing this identity.

\begin{enumerate}
\item The shared string $z$ has length $O(2^{-t}P(D+1))$ in every
history. The construction controls a quadratic energy of the wire
widths. Local tensor identities derived from a comparison of
Dirichlet energies reduce the leading coefficient in the energy bound
while raising the degree. The ports restore bounded degree at a fixed
cost. Grouping enough shrinking steps into each block gives the factor
$2^{-t}$ in the bound on total width.

\item The local contraction at any one $v$ uses at most $C^t$
occurrences of bit strings, map entries and original factors across
all its recursive levels, for a fixed $C>1$. Each string has length
at most $C^t(D+1)$. Here an original factor counts as one occurrence.
The maps use a fixed finite list of scalars satisfying fixed integer
polynomial equations. Applying the companion theorems over the prime
field gives a solution there, and the identities hold over any
commutative unital ring containing a solution. For given bit strings
a map entry is zero if a forbidden pattern occurs. Otherwise it is a
monomial in these scalars, of degree at most $C^t(D+1)$. The tests
and exponents take binary work polynomial in the string widths, with
an absolute exponent.

\item A history specifies the layout before any private bits are
read. Suppose $P,D\le V$ and $t=O(\log V)$. Given the initial arc
widths, indexed endpoint slots and masks, a nondeterministic
random-access procedure uses $O(P\operatorname{polylog}(2V))$ bit
operations and guessed bits per path. Its successful paths are in
bijection with admissible histories. At each level the procedure
records the fields, slices and ties needed for the local sums, and a
description of each map test of length $\operatorname{polylog}(2V)$.
These records can be accessed by address.
\end{enumerate}

\subsection{Evaluating closed networks}

The application to integer multiplication in \cite{oaT} uses both the cost
of evaluating the contraction and the cost of constructing its formula.

\begin{theorem}[Enumerable contraction]\label{thm:closed}
Fix a characteristic admitting ports. There is an $\eta>0$ such that
every closed bit network of volume at most $s$ and maximum degree at
most $s^{1/20}$, over any field of this characteristic, has a formula
for its contraction using $2^{O(s^{1-\eta})}$ arithmetic operations
and local entry queries. The topology determines the formula, which
can be generated deterministically in $2^{O(s^{1-\eta})}$ bit
operations and uses fixed prime-field constants. In characteristic
zero, for integer entries of polynomial bit length, evaluation takes
the same number of bit operations plus the time for the queries.
\end{theorem}

\begin{proof}
\emph{Routing the input network.}
For large $s$ choose the least $h$ with $P=h2^h\ge s^{9/10}$.
Then $P=\Theta(s^{9/10})$ and $h=\Theta(\log s)$. Use the positions
$(j,x)$ with $j\in\mathbb Z/h\mathbb Z$ and $x\in\{0,1\}^h$. From
$(j,x)$ put arcs to $(j+1,x)$ and $(j+1,x+e_j)$, where $e_j$ is the $j$th unit
vector. Add parallel slots of width zero to reach base degree $d$.

Place the input tensors in the $2^h$ lanes at $j=0$. Give each
tensor weight one plus its degree and always place the next tensor
in a lane of least current load. The total weight is at most $2s$.
The largest final load $q$ is at most the average load plus the
largest weight, so
\[
       q\le 2s/2^h+1+s^{1/20}=O(s^{1/10}\log s).
\]
Orient the original wires arbitrarily, and form a bipartite graph on
two copies of the lanes by joining the source lane of each wire to
its destination lane. Both sides have maximum degree at most $q$.
Pair their unused slots to complete a $q$-regular bipartite
multigraph. Successive perfect matchings split its edges into $q$
permutations.

Each permutation can be routed in one forward and one reverse trip
around the cursor $j$, by the Bene\v s construction~\cite{Benes}.
To see this, pair lanes differing in the first bit at both source
and destination. The permutation induces a degree-two bipartite
multigraph on the pairs. Color each cycle alternately. The switches
at the first and last stages can now send one color to each half, and
within either half the colored edges specify a permutation on one
fewer bit. Induction routes each permutation with at most one item
per lane on each pass. The reverse trip traverses the same arcs in
the opposite direction, which is allowed for bit equalities. Give
the two passes distinct bit slots and discard the padding wires.
Each original bit assignment has a unique extension satisfying the
equalities along these routes.
The maximum actual arc width $D$ and the local load, including the
tensors placed at $j=0$ and the routing ties, are both $O(q)$.

\emph{Enumerating the compressed expansion.}
Take $T=\lfloor a\log_2 s\rfloor$ blocks for a sufficiently small
fixed $a>0$ so that $b_0T\le h$. For each history the shared arc
bits have total length
\[
        O(2^{-T}P(D+1))=O(s^{1-a}\log s).
\]
The choices of local summands themselves use $O(PT\log s)$ bits.
There are $O(P)$ applications per substep and boundedly many substeps
per block. Each choice is a fixed number of integers of polynomial
magnitude. To enumerate histories deterministically, enumerate
the guess tapes for the layout procedure, checking choices in order
and requiring unused bits to be zero. This counts each successful
path once and costs $O(P\operatorname{polylog}s)$ bits per tape.

At each final position there are at most $C^T$ private fields of
length at most $C^T(D+1)$. Expanding the initial position factors
adds at most $O(C^T q)$ original tensors and routing ties. By taking
$a$ small enough, the number of private bits and the local load are
each at most $s^{1/2}$ for large $s$. For each pair $(\lambda,z)$
form the sum over $y_v$ at each position, then multiply the local
sums. For fixed $(\lambda,z)$ the work is the sum of the local enumeration
costs, because each private assignment occurs in a single local sum.
Thus the logarithm of the total work is at most
\[
 O\bigl(P\operatorname{polylog}s+s^{1-a}\log s+s^{1/2}\bigr)
       =O(s^{1-\eta})
\]
for some $\eta>0$. The layout records generate all indices, tests
and scalar monomials in polynomial work per assignment. This gives
both the arithmetic formula and the claimed bound for listing it.

In characteristic zero choose the fixed map scalars in $\mathbb Q$
and let $m$ clear their denominators. Each term of a local sum is
a product of polynomially many input entries and scalars with
polynomial total degree. Its denominator divides $m^{\mathrm{poly}(s)}$.
A local sum has at most $2^{\mathrm{poly}(s)}$ terms. Multiplying
polynomially many such sums and then adding over the outer enumeration
still gives polynomial bit length, since a common denominator at each
stage is a polynomial power of $m$. This argument also works for
inputs in a fixed ring of cyclotomic integers by using a fixed
integral basis. The polynomial cost per operation is absorbed into
the displayed bound.
\end{proof}
